\section{Introduction}\label{sec:intro}
Automated program repair (APR) has a long history in the programming languages and software engineering research communities. In the traditional setup, the APR system is given a bug-reproducing test suite and is tasked with fixing the bug. With the rise in machine learning methods, APR systems have increasingly relied on models (initially statistical and eventually deep learning-based) to perform critical APR tasks such as fault localization, patch generation, patch ranking and eventually end-to-end repair. More recently, work such as SWE-Agent~\cite{sweagent}, AutoCodeRover~\cite{autocoderover}, RepairAgent~\cite{repairagent},
CodeR~\cite{codeR},
AutoDev~\cite{autodev},
OpenDevin~\cite{opendevin}, and
others
have shown that, when incorporated into an agent-based system, LLMs can be used to perform end-to-end software engineering tasks in complex environments. Specifically, agentic repair systems can start from a bug description and autonomously generate bug-reproduction tests, localize faults, make candidate edits, validate these patches, and then submit a solution.

While such autonomous workflows have generated excitement in the APR community, systems in this space have been designed and evaluated using open-source bugs found in the GitHub ecosystem. In particular, SWE-Bench~\cite{swebench}, a collection of 2,294 Python bugs/fixes from popular GitHub repositories, and SWE-Bench-Lite, a subset of 300 bugs from SWE-Bench, have become the de-facto evaluation benchmarks for APR. %

It is not yet clear whether systems that perform well on SWE-Bench can achieve similar success when applied in the broader software industry, where we often encounter diverse collection of bugs which span a wide array of projects. Such conditions present both an opportunity and a challenge for APR systems. 
Given the enormous cost and effort to maintain code in enterprise environments, if agentic APR systems can perform as well in enterprise settings as on SWE-Bench, they hold the promise for substantial impact in industry.


To investigate the viability of agentic repair we first had to curate a benchmark set. Handling randomly-chosen bugs from Google's internal issue tracking system (GITS) would have been a non-starter for assessing agentic APR performance, due to various reasons, some of which are similar to those also encountered by the authors of SWE-Bench in their context: for example, each bug should have an easily-executable test.
Importantly, a randomly-chosen subset of bugs, while useful for population comparisons, provides little signal
for agent-level design and improvements.
Moreover, including bugs that
go beyond the current, but not future,
limitations of a basic agent
(e.g., screenshots) would
further cloud the informativeness
of any failures.

Consequently, we curated an evaluation set, GITS-Eval, of 178 bugs from Google's internal issue tracking system (GITS). 
This benchmark consists of 78 bugs reported by human developers and 100 machine-reported bugs, comprising 50 bugs reported by a suite of automated sanitizers (SAN), and 50 bugs reported by an automated test order dependency analyzer (TOD). These bugs reflect
different projects and programming languages, while remaining tractable for an APR system; the criteria for filtering are discussed in Section~\ref{sec:challengeset}.


To place GITS-Eval into context,  we also studied the differences
between SWE-Bench open source bugs and GITS. We sampled
2,000 bugs from GITS with filtering that reflects similar principles
to those underlying SWE-Bench and
found that GITS bugs exhibit different distributional characteristics from SWE-Bench open source bugs. In particular, we have observed differences in language diversity, size and spread of changes, and the presence of code terms in the bug description as a proxy for localization difficulty.
Specifically, we want to note that, due to these differences, the performance of an agent on one benchmark set may not be indicative of performance on the other.
In this paper, we make no claims on the performance of our baseline agentic system on SWE-Bench or SWE-Bench-Lite.

To establish a repair performance baseline on GITS-Eval, we built Passerine, a simple baseline agentic APR system inspired by SWE-Agent~\cite{sweagent} but designed to make use of Google’s internal development environment.
Using a limited command inventory, and with 20 sampled repair
runs (i.e. trajectories) per bug, Passerine is able to generate a
patch that passes bug tests (i.e. plausible) for 
25.6\% of the human-reported bugs, 73\% of machine-reported bugs (68\% of TOD-reported bugs and 78\% of SAN-reported bugs).
After manual annotation, we found that 17.9\% of human-reported bugs and 43\% of machine-reported bugs (24\% of TOD and 62\% of SAN) have at least one patch that was semantically equivalent to the ground-truth patch.
Thus, we establish that a simple agent can meaningfully solve issues from an industrially-relevant evaluation set. We expect
performance to improve as we
incorporate additional innovations
inspired by recent agentic
APR systems~\cite{specrover, autocoderover, codeR, marscode}.

We also make several observations about Passerine's behavior.
By analyzing command sequences, we note that Passerine adapts its behavior based
based on bug type (and associated bug information).
We found that Passerine can use rich bug reports, which carries implications for
bug report design in a future of increasingly prevalent agent-based systems.
Finally, we also found that trajectory analysis can reveal opportunities
for optimization, such as pruning degenerate trajectories.

















The remainder of the paper is organized as follows.
Section~\ref{sec:challengeset} presents our approach to understanding the differences between SWE-Bench and GITS bugs and shows how we construct a representative challenge set -- GITS-Eval -- appropriate for measuring APR system performance at Google.
Section~\ref{sec:passerine} presents Passerine, a SWE-Agent inspired agentic repair system designed to operate within Google's development environment.
Section~\ref{sec:results} presents our main results: an evaluation of Passerine's performance on a challenge set of 178 GITS bugs and broader observations
for agent-based APR in an enterprise context.
Section~\ref{sec:threats} describes threats and limitations
to our work.
Section~\ref{sec:relatedwork} discusses related work in 
multiple areas of APR.
